% Responsable: P5
\section{Comparación de enfoques}
\label{sec:comparacion_enfoques}

Los tres enfoques responden preguntas distintas. La CFG pregunta si la oración se puede derivar
con estas reglas; los modelos estadísticos (spaCy, Transformers), cuál es el análisis más
probable; y el LLM, qué respondería alguien a quien le dieron las reglas. La
tabla~\ref{tab:enfoques} resume cómo se comportó cada uno en el conjunto de prueba.

\begin{table}[h]
  \centering
  \small
  \begin{tabular}{p{0.2\linewidth}p{0.22\linewidth}p{0.24\linewidth}p{0.24\linewidth}}
    \toprule
    & CFG con rasgos & spaCy y Transformers & LLM con las reglas \\
    \midrule
    Decide la gramaticalidad & sí, exacto según las reglas & no (spaCy); según el español y la frecuencia (BETO) & sí, casi siempre según las reglas \\
    Dice por qué rechaza & sí, con la palabra y el rasgo & no & sí, aunque a veces se equivoca \\
    Fuera de la gramática & rechaza sin análisis & analiza casi todo & rechaza, pero etiqueta \\
    Concordancia & la detecta siempre & no la detecta & la detecta \\
    Ambigüedad & un solo árbol, por regla fija & un solo análisis, por probabilidad & un solo análisis; nunca dice \texttt{ambigua} \\
    Árboles & siempre válidos & otro formato (dependencias) & a veces fuera de la gramática \\
    Cómo se amplía & escribir reglas y palabras & reentrenar con un corpus anotado & cambiar el prompt \\
    \bottomrule
  \end{tabular}
  \caption{Comparación de los tres enfoques.}
  \label{tab:enfoques}
\end{table}

\paragraph{Robustez.} La CFG es frágil fuera de su fragmento. Una sola palabra fuera del léxico
(\emph{corre}) o una construcción sin regla (sujeto omitido, coordinación, subordinadas) hace que
rechace la oración sin dar ningún análisis. spaCy y Transformers degradan de forma gradual:
etiquetan cualquier oración, aunque se equivoquen más en las inusuales (\emph{*Perro el come}).
El LLM está en medio: con las reglas en el prompt rechaza lo que no se deriva, pero sigue
etiquetando bien todas las palabras, porque conoce el español.

\paragraph{Ambigüedad.} Ningún enfoque la maneja bien en la práctica. Nuestra CFG la detecta en
teoría (\emph{Él mató un elefante en sus pijamas} tiene dos árboles), pero el parser
descendente devuelve uno solo por una regla fija, que acierta 10 de 17 uniones de frases
preposicionales. Los modelos estadísticos eligen la lectura más probable y aciertan más (13 de
17 spaCy, 15 de 17 el LLM), pero tampoco avisan que hay otra. Para mostrar todas las lecturas
habría que usar un parser de gráfico sobre la gramática, que devuelve todos los árboles, y
elegir entre ellos con información de sentido, que es justo lo que la gramática no tiene.

\paragraph{Escalabilidad.} Ampliar la CFG es caro y lo hace una persona. Cada construcción nueva
necesita reglas, y cada regla puede crear conflictos nuevos: la palabra \emph{bajo} ya rompe el
análisis sin retroceso (sección~\ref{sec:parser}), y la coordinación o las subordinadas harían la
gramática mucho más ambigua. Escribir la concordancia dentro de la gramática multiplicaría las
reglas, por eso la verificamos aparte. Ampliar un modelo estadístico requiere un corpus anotado y
reentrenarlo, pero su cobertura ya es amplia. Ampliar el experimento con el LLM es lo más barato,
porque basta con cambiar el prompt, pero no da garantías: el modelo puede ignorar o aplicar mal
una regla nueva, como pasó con \emph{Leímos un libro}.

\paragraph{Recursos.} El parser es código Python y un léxico de 14~KB, y no depende de nada
externo. spaCy usa un modelo de 16~MB. Transformers necesita unos 1.5~GB para los dos modelos
y una biblioteca de aprendizaje profundo. El LLM no ocupa espacio local, pero depende de un
servicio externo: durante el experimento hubo errores por saturación del servicio y un límite
diario de llamadas que nos obligó a cambiar de modelo, y sus respuestas solo se pueden repetir
gracias al log.
